\documentclass[11pt]{article}
\usepackage[margin=2.4cm]{geometry}
\usepackage{fontspec}
\setmainfont{DejaVu Sans}
\usepackage{juliamono}
\usepackage{fvextra}
\usepackage{xcolor}
\usepackage{xurl}
\usepackage{hyperref}
\usepackage{amssymb}
\newcommand{\checkbox}{$\square$}
\newcommand{\answerchoices}{\checkbox\ YES \quad \checkbox\ NO \quad \checkbox\ CANNOT DETERMINE}
\newcommand{\decisionexplanation}{%
  \par\smallskip\textbf{If NO, give the reason for the mismatch (required). If
  CANNOT DETERMINE, state what context or evidence is missing:}\par
  \rule{\linewidth}{0.4pt}\par\vspace{7pt}\rule{\linewidth}{0.4pt}\par}
\setlength{\parindent}{0pt}\setlength{\parskip}{5pt}
\setlength{\emergencystretch}{2em}
\DefineVerbatimEnvironment{Vsrc}{Verbatim}{fontsize=\small,frame=single,
  breaklines=true,breakanywhere=true}
\DefineVerbatimEnvironment{Vlean}{Verbatim}{fontsize=\scriptsize,frame=single,
  numbers=left,numbersep=5pt,breaklines=true,breakanywhere=true}
\hypersetup{colorlinks=true,urlcolor=blue}

\begin{document}
\begin{center}{\LARGE Proof Verification Document}\\[2pt]{\large Algorithms\_27\_\_sol3}\\[2pt]{\small QEDBench, second-model flags (confirmed)}\end{center}
\begin{center}\fcolorbox{blue!60!black}{blue!4}{\parbox{0.92\linewidth}{\textbf{Make three independent decisions.} Follow the evidence chain: original source $\longrightarrow$ challenged claim $\longrightarrow$ proposed counterexample $\longrightarrow$ Lean theorem. A record is valid only when all three links are correct. Complete the human sign-off before reading the automated review appendix; those earlier results are background, not an answer key.}}\end{center}
\textbf{Use only this document, except for the public source link.} Everything needed for Questions 2 and 3 is included here. For Question 1, open the public source link when one is shown. If there is no public link, the complete stored problem and solution are included here. If the information is still not enough, choose CANNOT DETERMINE.
For each question, choose \textbf{YES}, \textbf{NO}, or \textbf{CANNOT DETERMINE}. Choose CANNOT DETERMINE if the document does not provide enough information for you to decide; do not mark YES if you are uncertain. If you find a specific mismatch, choose NO and explain it in the space below that question. Do not rewrite the claim, add assumptions, or substitute a different counterexample.
\textbf{How copied text is shown.} The problem, quoted passage, and proposed counterexample come from the stored record. The Lean theorem statement and Lean files are copied exactly from the certificate. Some long lines continue on the next printed line; this is only page formatting and does not change the copied text. Line numbers are added only to help you find passages. If the stored record contains a character that cannot be printed, the document shows \texttt{[U+XXXX CONTROL]} in its place.
\begin{description}
\item[Challenged claim.] The exact statement from the source that this record says is false.
\item[Proposed counterexample.] A concrete case that should satisfy the claim's assumptions but make its conclusion false.
\item[Lean certificate.] A Lean file proving a formal theorem. Successful compilation proves the theorem as written, not that it is the right translation of the source.
\end{description}
\section*{Question 1 --- Does this proof verification document accurately represent the source?}
\textbf{What to decide.} Compare the challenged passage below with the original source---either the public link or the complete source text in this document. Answer YES only if they make the same claim and include the same assumptions or conditions. If this document changes the meaning, leaves out a condition, or points to a different passage, answer NO.
\subsection*{Original source}
\textbf{Source record:}
\begin{Vsrc}
qedbench_certificate_availability_triage/data/records_alg_graph.jsonl
\end{Vsrc}
\textit{No public web page is recorded for this item. The complete stored problem and solution are included below, so you do not need to open the source-record file.}
\subsection*{Problem or question from the source (verbatim)}
\begin{Vsrc}
Problem Statement
Problem: Algorithms - 27
Consider the decision problem denoted as **12Partition**. The input consists of a set
S containing 12n positive integers. The problem requires determining if there exists a
partition of S into n disjoint subsets A1, A2, . . . , An such that the cardinality of each
subset |Ai| is 12, and the sum of the elements within each subset is identical. Provide a
proof establishing that **12Partition** is NP-hard.
1
\end{Vsrc}
\subsection*{Complete stored solution (verbatim)}
\begin{Vsrc}
Solution:
Consider the decision problem denoted as 12Partition. The input consists of
a set S containing 12n positive integers. The problem requires determining if there exists a
partition of S into n disjoint subsets A1, A2, . . . , An such that the cardinality of each subset
|Ai| is 12, and the sum of the elements within each subset is identical. Provide a proof
establishing that 12Partition is NP-hard.
Theorem 1. The 12Partition problem is NP-hard.
Proof. We prove NP-hardness by a polynomial-time reduction from the 3-Partition prob-
lem, which is known to be strongly NP-complete [?].
Input: A set T = {t1, t2, . . . , t3m} of 3m positive integers and a positive integer B such
that:
• P3m
i=1 ti = mB
•
B
4 < ti < B
2 for all i
Question: Can T be partitioned into m disjoint subsets T1, T2, . . . , Tm such that |Tj| = 3
and P
t∈Tj t = B for all j?
Given an instance (T, B) of 3-Partition with |T| = 3m, we construct an instance S of
12Partition as follows:
For each element ti ∈T (where i = 1, 2, . . . , 3m), we create four distinct elements in S:
s(1)
i
= ti+K·43m+1−i,
s(2)
i
= ti+2K·43m+1−i,
s(3)
i
= ti+3K·43m+1−i,
s(4)
i
= ti+4K·43m+1−i
(1)
where K = 1+P3m
i=1 ti = 1+mB is chosen to be larger than the sum of all original elements.
The resulting set S contains 12m elements. Note that this construction can be computed
in polynomial time.
Lemma 1. For any i ̸= i′, the positional weights 43m+1−i and 43m+1−i′ differ by at least a
factor of 4.
Lemma 2. For any fixed index i, we have:
4
X
j=1
s(j)
i
= 4ti + 10K · 43m+1−i
(2)
Lemma 3. In any valid 12Partition of S, each subset of 12 elements must contain exactly
all four copies s(1)
i , s(2)
i , s(3)
i , s(4)
i
for exactly three distinct indices i.
Proof of Lemma ??. Consider any subset A in a valid 12Partition.
Since |A| = 12 and
|S| = 12m, there are m such subsets total.
For each index i ∈{1, . . . , 3m}, let ni denote the number of copies of the form s(j)
i
contained in A. Clearly, 0 ≤ni ≤4 and P3m
i=1 ni = 12.
The sum of elements in A can be written as:
X
s∈A
s =
X
i:ni>0
X
j∈Ji
s(j)
i
!
(3)
1

where Ji ⊆{1, 2, 3, 4} with |Ji| = ni.
For the partition to be valid, all subsets must have the same sum. Since the total sum
of all elements in S is:
3m
X
i=1
4
X
j=1
s(j)
i
= 4
3m
X
i=1
ti + 10K
3m
X
i=1
43m+1−i = 4mB + 10K
3m
X
i=1
43m+1−i
(4)
Each of the m subsets must have sum:
Σ = 4B + 10K · 1
m
3m
X
i=1
43m+1−i
(5)
Now, consider the contribution from the K · 43m+1−i terms. Since these are powers of 4
with different exponents, and K is chosen large enough, the sum modulo K must balance
correctly.
Specifically, if we don’t take all four copies for a given index i (i.e., 0 < ni < 4), then the
coefficient of K · 43m+1−i in our subset sum would be (1 + 2 + · · · + k) for some k < 4, which
is at most 1 + 2 + 3 = 6. But we need the coefficient to be 10 (on average) for each of three
indices.
By a careful accounting argument using the distinctness of powers of 4 and the constraint
that P
i ni = 12, we can show that the only way to achieve the required sum is if ni ∈{0, 4}
for all i, with exactly three indices having ni = 4 and all others having ni = 0.
Forward Direction: Suppose (T, B) has a valid 3-partition {T1, . . . , Tm} where each
Tj = {ti1, ti2, ti3} with ti1 + ti2 + ti3 = B.
For each Tj, define Aj to be the set containing all four copies of each of the three elements
in Tj:
Aj = {s(ℓ)
ik : k ∈{1, 2, 3}, ℓ∈{1, 2, 3, 4}}
(6)
Then |Aj| = 12 and by Lemma ??:
X
s∈Aj
s = 4(ti1 + ti2 + ti3) + 10K(43m+1−i1 + 43m+1−i2 + 43m+1−i3) = 4B + 10K · Cj
(7)
where Cj depends only on the indices. Since each index appears in exactly one partition set,
all Aj have the same sum.
Reverse Direction: Suppose S has a valid 12Partition {A1, . . . , Am}. By Lemma ??,
each Aj contains all four copies of exactly three distinct indices from T. Let these indices
for Aj be {ij
1, ij
2, ij
3}.
Define Tj = {tij
1, tij
2, tij
3}. Since all Aj have the same sum and the positional weights
balance out symmetrically, we must have:
tij
1 + tij
2 + tij
3 = B for all j
(8)
Thus {T1, . . . , Tm} is a valid 3-partition of T.
Since 3-Partition is NP-complete and our reduction is polynomial-time computable, we
conclude that 12Partition is NP-hard.
2
\end{Vsrc}
\subsection*{Challenged passage recorded from the source (stored text)}
\begin{Vsrc}
In any valid 12Partition of S, each subset of 12 elements must contain exactly
\end{Vsrc}
\subsection*{Surrounding source lines (stored text with added line labels)}
The labels refer to lines in the solution text stored with the source record; the source-line content is not transliterated.
\begin{Vsrc}
0003 | a set S containing 12n positive integers. The problem requires determining if there exists a
0004 | partition of S into n disjoint subsets A1, A2, . . . , An such that the cardinality of each subset
0005 | |Ai| is 12, and the sum of the elements within each subset is identical. Provide a proof
0006 | establishing that 12Partition is NP-hard.
0007 | Theorem 1. The 12Partition problem is NP-hard.
0008 | Proof. We prove NP-hardness by a polynomial-time reduction from the 3-Partition prob-
0009 | lem, which is known to be strongly NP-complete [?].
0010 | Input: A set T = {t1, t2, . . . , t3m} of 3m positive integers and a positive integer B such
0011 | that:
0012 | • P3m
0013 | i=1 ti = mB
0014 | •
0015 | B
0016 | 4 < ti < B
0017 | 2 for all i
0018 | Question: Can T be partitioned into m disjoint subsets T1, T2, . . . , Tm such that |Tj| = 3
0019 | and P
0020 | t∈Tj t = B for all j?
0021 | Given an instance (T, B) of 3-Partition with |T| = 3m, we construct an instance S of
0022 | 12Partition as follows:
0023 | For each element ti ∈T (where i = 1, 2, . . . , 3m), we create four distinct elements in S:
0024 | s(1)
0025 | i
0026 | = ti+K·43m+1−i,
0027 | s(2)
\end{Vsrc}
\textbf{A YES means:} the proof verification document faithfully represents both the source context and the challenged claim, without dropping a condition or changing its meaning.
\section*{Question 2 --- Does the proposed counterexample disprove that claim?}
\textbf{Claim to test (repeated for convenience).}
\begin{Vsrc}
In any valid 12Partition of S, each subset of 12 elements must contain exactly
\end{Vsrc}
\subsection*{How the record says the argument uses this claim}
This is the record's description, not an established fact; check it against the source.
\begin{Vsrc}
Lemma 3 (the load-bearing structural lemma of the reduction): "In any valid 12Partition of S, each subset of 12 elements must contain exactly all four copies s(1)i, s(2)i, s(3)i, s(4)i for exactly three distinct indices i." A single explicit instance refutes this lemma and, with it, the correctness of the reduction: a NO instance of 3-Partition is mapped to a YES instance of 12Partition.
\end{Vsrc}
\subsection*{Proposed counterexample (exact values or construction)}
\begin{Vsrc}
Take the 3-Partition instance A = (5, 5, 5, 5, 6, 8) with m = 2, B = 17. It is legitimate: sum(A) = 34 = mB, and every element satisfies B/4 < a_i < B/2 (checked exactly as 4*a_i > B and 2*a_i < B: 20>17, 24>17, 32>17 and 10<17, 12<17, 16<17). A is a NO instance of 3-Partition: the only triple sums available are 5+5+5=15, 5+5+6=16, 5+5+8=18, 5+6+8=19, so no triple sums to B = 17 at all (exhaustive check over all 2-triple partitions).

The reduction builds S = four labelled copies of each element: sixteen 5's (copies of indices 1-4), four 6's (index 5), four 8's (index 6); |S| = 24, sum(S) = 136 = 4mB, target block sum 4B = 68, block size 12. S HAS a valid 12Partition:

  Block 1 = { s^(1)_1 ; s^(1)_3, s^(2)_3, s^(3)_3 ; s^(1)_4, s^(2)_4, s^(3)_4, s^(4)_4 ; s^(1)_5, s^(2)_5 ; s^(1)_6, s^(2)_6 }
          values [5,5,5,5,5,5,5,5,6,6,8,8], size 12, sum = 8*5 + 2*6 + 2*8 = 68 = 4B
  Block 2 = { s^(2)_1, s^(3)_1, s^(4)_1 ; s^(1)_2, s^(2)_2, s^(3)_2, s^(4)_2 ; s^(4)_3 ; s^(3)_5, s^(4)_5 ; s^(3)_6, s^(4)_6 }
          values [5,5,5,5,5,5,5,5,6,6,8,8], size 12, sum = 68 = 4B

The two blocks are disjoint and together contain each of the 24 labelled copies exactly once (verified independently by explicit set comparison).

This directly contradicts Lemma 3: block 1 contains all four copies only of index 4, and takes PARTIAL copy sets of indices 1 (one copy), 3 (three copies), 5 (two copies), 6 (two copies) — it is not "all four copies for exactly three distinct indices". Same for block 2. And it breaks the reduction: a NO instance of 3-Partition is mapped to a YES instance of 12Partition.

Two further witnesses printed by the same script:
 - Smallest found overall: A = (4,4,4,4,4,6), m = 2, B = 13 (triple sums only 12 and 14, so NO), yet S splits into two blocks of ten 4's + two 6's = 52 = 4B each.
 - m = 3: A = (4,4,4,4,4,4,4,5,6), B = 13, a NO instance (no partition into three triples of sum 13), yet S splits into three valid 12-blocks of sum 52.
\end{Vsrc}
\subsection*{Why the record says the counterexample works}
Check this explanation independently; it is evidence to assess, not an answer key.
\begin{Vsrc}
Lemma 3 is a finitely refutable statement about finite integer instances, so a single explicit instance suffices. The failure is structural, not accidental: in 3-Partition the bound B/4 < a_i < B/2 forces exactly three elements per group of sum B. After quadrupling, the 12-element blocks must sum to 4B while each element lies in (4B/16, 4B/8) = (T/16, T/8) with T = 4B — a window WIDER than the (T/24, T/12) needed to force the "4 copies of 3 indices" structure. A block of 12 such elements can have any sum strictly between 0.75T and 1.5T, so T = 4B is reachable by many multiplicity patterns with 0 < c_i < 4, not just by c_i in {0,4}. The lemma silently assumes the size window transfers under quadrupling; it does not.

Concretely with A = (5,5,5,5,6,8), B = 17: the intended (and only lemma-permitted) block sums are 4*(a_i+a_j+a_k), which take values 60, 64, 72, 76 — never 68 = 4B, exactly because A is a NO instance. But the mixed pattern 8 fives + 2 sixes + 2 eights = 68 hits the target while using 1, 3, 4, 2, 2 copies of the various indices, and its complement is the identical pattern, so both blocks are valid. Hence a valid 12Partition exists with no block of the lemma's shape.

The consequence is that the claimed reduction is not sound: soundness needs "12Partition YES => 3-Partition YES", and here the 3-Partition instance is provably NO (verified by exhaustive enumeration of all partitions into triples) while its image is provably YES (exhibited partition verified element-by-element on labelled copies). The counterexample is not an artifact of equal values alone — the m=2, B=17 instance has three distinct values, and an m=3 instance with three distinct values was also found — so it cannot be repaired by tie-breaking or perturbation arguments.

All arithmetic in the search and in the independent re-verification is exact integer arithmetic (no floats anywhere; the size condition is tested as 4*a > B and 2*a < B), the enumeration order is fixed, and the script is deterministic and stdlib-only. Runtime 0.6 s.
\end{Vsrc}
\subsection*{How to decide}
\begin{enumerate}
\item List every assumption that applies to the challenged claim.
\item Check that the proposed example satisfies every one of those assumptions.
\item Recompute the claimed outcome and the actual outcome using exact reasoning.
\item Confirm that the conflict is with the challenged claim itself, not an unrelated sentence.
\end{enumerate}
\textbf{A YES means:} the example is fully specified, satisfies all relevant assumptions, and makes the challenged claim false. If a necessary value, condition, or calculation is missing, choose CANNOT DETERMINE.
\section*{Question 3 --- Does the final Lean theorem statement verify this same counterexample?}
\textbf{What to decide.} A Lean file can compile while proving the wrong theorem. Compare these three links independently:
\begin{enumerate}
\item The Lean definitions represent the same mathematical objects, assumptions, and operations as the challenged claim.
\item The Lean construction uses the same concrete values or example shown in Question 2.
\item The final theorem statement asserts that the challenged claim fails on that example, not merely a weaker or nearby statement.
\end{enumerate}
\subsection*{Selected final Lean theorem statement (verbatim UTF-8 excerpt)}
This is an exact contiguous slice of the certificate file, from the theorem keyword through the original \texttt{:=} token where its proof begins. It is shown first so the formal target is not buried in the complete file.
\begin{Vlean}
theorem Algorithms_27__sol3_refutation :
    legit Aw 2 17 = true
  ∧ hasTP 17 2 Aw = false
  ∧ validPartition Aw 17 2 [block1, block2] = true
  ∧ (List.range 6).map (copiesOf block1) = [1, 0, 3, 4, 2, 2]
  ∧ (List.range 6).map (copiesOf block2) = [3, 4, 1, 0, 2, 2]
  ∧ lemmaShape 6 block1 = false
  ∧ lemmaShape 6 block2 = false
  ∧ lemma3At Aw 17 2 [block1, block2] = false
  ∧ soundnessAt Aw 17 2 [block1, block2] = false :=
\end{Vlean}
\subsection*{Project-local imported Lean files (complete, verbatim UTF-8)}
This section follows the certificate's import commands recursively and prints every repository-owned Lean source file on which the final theorem statement depends. Each file is copied in full, including its definitions and controls. External Mathlib, Std, and Lean Init source is not duplicated.
No project-local imported Lean files were found. External Mathlib, Std, and Lean Init source is not duplicated in this proof verification document.
\subsection*{Complete counterexample Lean file (verbatim UTF-8)}
The complete Lean file is included below so you can check every definition and supporting result used by the final theorem statement. Use it together with the project-local imported Lean files above.
\textbf{Lean certificate file:}
\begin{Vsrc}
partB_feasibility/qb_opus_certificates/Algorithms_27__sol3.lean
\end{Vsrc}
\begin{Vlean}
/-
================================================================================
  RECORD:  Algorithms_27__sol3   (Domain: Algorithms — QEDBench templated proof)
  SOURCE:  QEDBench bank, record Algorithms_27, solution 3 (no public URL)
  ROUTE:   SELF_CONTAINED_LEAN   (no imports, no axioms, no `sorry`, no
                                  `native_decide`; every theorem is closed by
                                  kernel-checked `decide` on 6-element instances
                                  with 24 labelled copies)
           The V2 DSL was checked first: QEDTemplateV2 has NO constructor for
           multiset partition / 3-Partition / block-sum claims (its only
           "partition" claims concern max-cut vertex partitions), so the claim
           is not representable there and this certificate is self-contained.
================================================================================

CLAIM BEING REFUTED (Lemma 3 of the 3-Partition -> 12Partition reduction):

    "In any valid 12Partition of S, each subset of 12 elements must contain
     exactly all four copies s^(1)_i .. s^(4)_i for exactly three distinct
     indices i."

Here A = (a_1, ..., a_3m) is a legitimate 3-Partition instance
(sum(A) = mB and B/4 < a_i < B/2 for every i), S consists of FOUR labelled
copies of each a_i (|S| = 12m, sum(S) = 4mB), and a valid 12Partition of S
is a partition of the 12m labelled copies into m subsets of exactly 12
elements, each of sum 4B.  Lemma 3 is what makes the claimed reduction
sound ("12Partition YES => 3-Partition YES"); both the lemma and the
soundness direction are universally quantified over finite integer
instances, hence finitely refutable by a single explicit instance.

WITNESS (primary; verified independently by
partB_feasibility/counterexample_search_qb_opus/verify_Algorithms_27__sol3.py):

    A = (5, 5, 5, 5, 6, 8), m = 2, B = 17.
    Legitimacy: sum(A) = 34 = 2*17, and the window holds exactly:
    4*a_i > 17 (20, 20, 20, 20, 24, 32) and 2*a_i < 17 (10, 10, 10, 10, 12, 16).
    A is a NO instance of 3-Partition: the available triple sums are
    5+5+5 = 15, 5+5+6 = 16, 5+5+8 = 18, 5+6+8 = 19 — never 17
    (decided below by a complete anchored backtracking search).

    Yet S = { s^(c)_i : 1 <= i <= 6, 1 <= c <= 4 } HAS a valid 12Partition
    (0-based below: element (i, c) is copy c of index i, value A[i]):

      block1 = (0,0); (2,0),(2,1),(2,2); (3,0),(3,1),(3,2),(3,3);
               (4,0),(4,1); (5,0),(5,1)
               values [5,5,5,5,5,5,5,5,6,6,8,8], size 12, sum 68 = 4B
      block2 = (0,1),(0,2),(0,3); (1,0),(1,1),(1,2),(1,3); (2,3);
               (4,2),(4,3); (5,2),(5,3)
               values [5,5,5,5,5,5,5,5,6,6,8,8], size 12, sum 68 = 4B

    The blocks are disjoint and cover all 24 labelled copies exactly once.
    Per-index copy multiplicities: block1 uses (1, 0, 3, 4, 2, 2) and
    block2 uses (3, 4, 1, 0, 2, 2) — block1 contains all four copies only
    of index 3 and PARTIAL copy sets of indices 0, 2, 4, 5, so it is not
    "all four copies for exactly three distinct indices"; same for block2.
    Lemma 3 is false, and the reduction maps a NO instance of 3-Partition
    to a YES instance of 12Partition, so it is unsound.

SECONDARY WITNESS (smallest found): A = (4, 4, 4, 4, 4, 6), m = 2, B = 13
(triple sums only 12 and 14, so NO), yet S splits into two valid blocks of
ten 4's + two 6's = 52 = 4B each.

ROOT CAUSE:  In 3-Partition the window B/4 < a_i < B/2 forces exactly three
elements per group of sum B.  After quadrupling, a 12-element block must sum
to T = 4B while each element lies in (T/16, T/8) — a window WIDER than the
(T/24, T/12) needed to force the "4 copies of 3 indices" structure, so T is
reachable by mixed multiplicity patterns with copy counts strictly between
0 and 4.  The lemma silently assumes the size window transfers under
quadrupling; it does not.

MODELLING JUDGMENTS (each pinned by a control theorem below):

 1. VALUES AND LABELS.  An instance is a `List Nat` of values (0-based
    index i carries value `nthVal vals i`); a labelled copy is a pair
    `(i, c) : Nat × Nat` with `c < 4`.  `labelled n` enumerates the full
    labelled set S; its size (24) and duplicate-freeness are pinned by
    `Algorithms_27__sol3_checker_discriminates`.  Nat values suffice: the
    witness values are positive integers, and a refutation for positive
    integers refutes the general claim a fortiori.

 2. LEGITIMACY.  `legit vals m B` checks |A| = 3m, sum(A) = mB, and the
    window EXACTLY as integers: `B < 4*a` and `2*a < B` (no division, no
    floats) — precisely the reduction's stated precondition.

 3. 3-PARTITION DECIDED, NOT ASSUMED.  `hasTP B fuel l` is a COMPLETE
    anchored backtracking search: the first remaining element anchors a
    triple, all (ordered) pairs of the rest complete it, recursion on the
    remainder with fuel = number of triples.  Any partition into triples of
    sum B can be reordered to put its first triple's members first, so
    anchoring loses nothing; ordered pairs only duplicate work.  Fuel m
    suffices (each level removes one triple from 3m elements).  Control:
    `hasTP` returns `true` on the legitimate YES instance (5,5,7,5,6,6)
    (5+5+7 = 17, 5+6+6 = 17), so the `false` verdict on the witness is the
    output of a genuinely complete decision procedure, not a stub.

 4. VALID 12PARTITION.  `validPartition vals B m blocks` demands: exactly m
    blocks; each block has exactly 12 labelled copies summing to 4B; the
    concatenation of the blocks is duplicate-free and equals the labelled
    set S as a set (mutual containment; both sides duplicate-free).  This
    is exactly the definition quoted in the claim.  Controls show it
    rejects a duplicated block list and a wrong block count.

 5. LEMMA 3 SHAPE.  `lemmaShape n b` holds iff every index contributes 0 or
    4 copies to b and exactly three indices contribute 4.  (For a 12-element
    block this is exactly "contains all four copies of exactly three
    distinct indices".)  Control: it ACCEPTS the intended all-copies blocks
    of the YES instance, so its rejection of block1/block2 is not vacuous.

 6. THE CLAIM AS AN INSTANCE.  `lemma3At vals B m blocks` is the exact
    instance "blocks form a valid 12Partition => every block has the lemma
    shape"; `soundnessAt vals B m blocks` is the exact instance
    "(legitimate A) and (blocks form a valid 12Partition of S) => A is a
    3-Partition YES instance".  Both are `false` on the witness and `true`
    on the intended partition of the YES control instance.

 7. This mirrors the exhaustive search and independent re-verification in
    partB_feasibility/counterexample_search_qb_opus/Algorithms_27__sol3.py
    and verify_Algorithms_27__sol3.py (exact integer arithmetic throughout).
================================================================================
-/

set_option maxRecDepth 10000
set_option maxHeartbeats 4000000

namespace Algorithms27sol3

/-! ### Generic machinery -/

/-- Sum of a `Nat` list. -/
def sumL : List Nat → Nat
  | [] => 0
  | x :: xs => x + sumL xs

/-- `i`-th value of the instance (0 out of range — never load-bearing: every
label used below is in range, pinned by the set-equality with `labelled`). -/
def nthVal : List Nat → Nat → Nat
  | [], _ => 0
  | x :: _, 0 => x
  | _ :: xs, i + 1 => nthVal xs i

/-- Flatten a list of lists. -/
def joinL {α : Type} : List (List α) → List α
  | [] => []
  | l :: ls => l ++ joinL ls

/-- Equality of labelled copies. -/
def pairEq (p q : Nat × Nat) : Bool := p.1 == q.1 && p.2 == q.2

/-- Membership of a labelled copy in a list. -/
def memP (l : List (Nat × Nat)) (p : Nat × Nat) : Bool := l.any (pairEq p)

/-- No labelled copy listed twice. -/
def noDupP : List (Nat × Nat) → Bool
  | [] => true
  | p :: ps => !(memP ps p) && noDupP ps

/-! ### Legitimacy of a 3-Partition instance -/

/-- `vals` is a legitimate 3-Partition instance for `(m, B)`:
`|A| = 3m`, `sum(A) = mB`, and the exact integer window
`B < 4*a` and `2*a < B` (i.e. `B/4 < a < B/2` over ℚ) for every element. -/
def legit (vals : List Nat) (m B : Nat) : Bool :=
  (vals.length == 3 * m) && (sumL vals == m * B)
  && vals.all (fun a => decide (B < 4 * a) && decide (2 * a < B))

/-! ### Complete 3-Partition decision (anchored backtracking) -/

/-- All ways to pick one element of a list, with the remaining list. -/
def oneWithRest : List Nat → List (Nat × List Nat)
  | [] => []
  | x :: xs => (x, xs) :: (oneWithRest xs).map (fun yr => (yr.1, x :: yr.2))

/-- All ways to pick an (ordered) pair of elements, with the remaining list. -/
def pairsWithRest (l : List Nat) : List (Nat × Nat × List Nat) :=
  joinL ((oneWithRest l).map (fun ar =>
    (oneWithRest ar.2).map (fun br => (ar.1, br.1, br.2))))

/-- `hasTP B fuel l`: `l` can be partitioned into triples each of sum `B`
(complete search: the head anchors a triple, every pair of the rest is tried,
recursion on the remainder; `fuel` = number of triples = structural measure). -/
def hasTP (B : Nat) : Nat → List Nat → Bool
  | _, [] => true
  | 0, _ :: _ => false
  | fuel + 1, x :: xs =>
      (pairsWithRest xs).any (fun t =>
        (x + t.1 + t.2.1 == B) && hasTP B fuel t.2.2)

/-! ### Labelled copies and valid 12Partitions -/

/-- The labelled set S: four copies `(i, 0), ..., (i, 3)` of every index
`i < n` (for the witness, n = 6: 24 labelled copies). -/
def labelled (n : Nat) : List (Nat × Nat) :=
  joinL ((List.range n).map (fun i => (List.range 4).map (fun c => (i, c))))

/-- Value carried by a block of labelled copies. -/
def blockSum (vals : List Nat) (b : List (Nat × Nat)) : Nat :=
  sumL (b.map (fun p => nthVal vals p.1))

/-- `blocks` is a VALID 12PARTITION of S: exactly `m` blocks, each of exactly
12 labelled copies summing to `4B`, jointly duplicate-free and equal AS A SET
to the labelled set S (mutual containment; both sides duplicate-free). -/
def validPartition (vals : List Nat) (B m : Nat)
    (blocks : List (List (Nat × Nat))) : Bool :=
  let S := labelled vals.length
  let allC := joinL blocks
  (blocks.length == m)
  && blocks.all (fun b => (b.length == 12) && (blockSum vals b == 4 * B))
  && noDupP allC
  && allC.all (fun p => memP S p)
  && S.all (fun p => memP allC p)

/-- Number of copies of index `i` inside block `b`. -/
def copiesOf (b : List (Nat × Nat)) (i : Nat) : Nat :=
  (b.filter (fun p => p.1 == i)).length

/-- Lemma 3's shape: `b` contains all four copies of exactly three distinct
indices (every index contributes 0 or 4 copies; exactly three contribute 4). -/
def lemmaShape (n : Nat) (b : List (Nat × Nat)) : Bool :=
  (List.range n).all (fun i => (copiesOf b i == 0) || (copiesOf b i == 4))
  && (((List.range n).filter (fun i => copiesOf b i == 4)).length == 3)

/-- LEMMA 3, on one instance: if `blocks` is a valid 12Partition of S then
EVERY block has the lemma shape. -/
def lemma3At (vals : List Nat) (B m : Nat)
    (blocks : List (List (Nat × Nat))) : Bool :=
  !(validPartition vals B m blocks)
  || blocks.all (fun b => lemmaShape vals.length b)

/-- SOUNDNESS of the reduction, on one instance: if `vals` is legitimate and
`blocks` is a valid 12Partition of S, then `vals` is a 3-Partition YES
instance. -/
def soundnessAt (vals : List Nat) (B m : Nat)
    (blocks : List (List (Nat × Nat))) : Bool :=
  !(legit vals m B && validPartition vals B m blocks) || hasTP B m vals

/-! ### The witness -/

/-- A = (5, 5, 5, 5, 6, 8), m = 2, B = 17. -/
def Aw : List Nat := [5, 5, 5, 5, 6, 8]

/-- Block 1 of the exhibited 12Partition (labelled copies, 0-based). -/
def block1 : List (Nat × Nat) :=
  [(0, 0),
   (2, 0), (2, 1), (2, 2),
   (3, 0), (3, 1), (3, 2), (3, 3),
   (4, 0), (4, 1),
   (5, 0), (5, 1)]

/-- Block 2 of the exhibited 12Partition. -/
def block2 : List (Nat × Nat) :=
  [(0, 1), (0, 2), (0, 3),
   (1, 0), (1, 1), (1, 2), (1, 3),
   (2, 3),
   (4, 2), (4, 3),
   (5, 2), (5, 3)]

/-! ### Secondary witness (smallest found) -/

/-- A = (4, 4, 4, 4, 4, 6), m = 2, B = 13. -/
def A2w : List Nat := [4, 4, 4, 4, 4, 6]

def blockS1 : List (Nat × Nat) :=
  [(0, 0), (0, 1), (0, 2), (0, 3),
   (1, 0), (1, 1), (1, 2), (1, 3),
   (2, 0), (2, 1),
   (5, 0), (5, 1)]

def blockS2 : List (Nat × Nat) :=
  [(2, 2), (2, 3),
   (3, 0), (3, 1), (3, 2), (3, 3),
   (4, 0), (4, 1), (4, 2), (4, 3),
   (5, 2), (5, 3)]

/-! ### Control instance (legitimate YES instance, intended-shape partition) -/

/-- A = (5, 5, 7, 5, 6, 6), m = 2, B = 17: a legitimate YES instance
(5 + 5 + 7 = 17 and 5 + 6 + 6 = 17). -/
def Ayes : List Nat := [5, 5, 7, 5, 6, 6]

/-- The reduction's INTENDED block for the triple {0, 1, 2}: all copies. -/
def confBlock1 : List (Nat × Nat) :=
  [(0, 0), (0, 1), (0, 2), (0, 3),
   (1, 0), (1, 1), (1, 2), (1, 3),
   (2, 0), (2, 1), (2, 2), (2, 3)]

/-- The intended block for the triple {3, 4, 5}: all copies. -/
def confBlock2 : List (Nat × Nat) :=
  [(3, 0), (3, 1), (3, 2), (3, 3),
   (4, 0), (4, 1), (4, 2), (4, 3),
   (5, 0), (5, 1), (5, 2), (5, 3)]

/-! ================================================================
    MAIN REFUTATION
    ================================================================ -/

/-- **Refutation of Lemma 3 and of the reduction's soundness.**
A = (5,5,5,5,6,8), m = 2, B = 17 is a LEGITIMATE 3-Partition instance
(sum = mB, exact window B/4 < a < B/2) that is a NO instance (complete
backtracking finds no partition into triples of sum 17).  Yet
[block1, block2] is a VALID 12Partition of its labelled set S — 2 blocks
of 12 labelled copies each, each summing to 4B = 68, jointly covering all
24 copies exactly once — and NEITHER block has Lemma 3's shape: the pinned
per-index multiplicities are (1, 0, 3, 4, 2, 2) and (3, 4, 1, 0, 2, 2),
i.e. partial copy sets, not "all four copies of exactly three indices".
Hence `lemma3At = false` (Lemma 3 fails on a valid 12Partition) and
`soundnessAt = false` (a NO instance maps to a YES instance). -/
theorem Algorithms_27__sol3_refutation :
    legit Aw 2 17 = true
  ∧ hasTP 17 2 Aw = false
  ∧ validPartition Aw 17 2 [block1, block2] = true
  ∧ (List.range 6).map (copiesOf block1) = [1, 0, 3, 4, 2, 2]
  ∧ (List.range 6).map (copiesOf block2) = [3, 4, 1, 0, 2, 2]
  ∧ lemmaShape 6 block1 = false
  ∧ lemmaShape 6 block2 = false
  ∧ lemma3At Aw 17 2 [block1, block2] = false
  ∧ soundnessAt Aw 17 2 [block1, block2] = false := by
  decide

/-- **Secondary witness (smallest found).**  A = (4,4,4,4,4,6), m = 2,
B = 13: legitimate, a NO instance (triple sums are only 12 and 14), yet its
labelled set splits into two valid blocks of ten 4's + two 6's = 52 = 4B,
again with partial copy sets (multiplicities (4,4,2,0,0,2) / (0,0,2,4,4,2)).
The failure is therefore not special to the primary witness. -/
theorem Algorithms_27__sol3_secondary :
    legit A2w 2 13 = true
  ∧ hasTP 13 2 A2w = false
  ∧ validPartition A2w 13 2 [blockS1, blockS2] = true
  ∧ (List.range 6).map (copiesOf blockS1) = [4, 4, 2, 0, 0, 2]
  ∧ lemmaShape 6 blockS1 = false
  ∧ lemmaShape 6 blockS2 = false
  ∧ lemma3At A2w 13 2 [blockS1, blockS2] = false
  ∧ soundnessAt A2w 13 2 [blockS1, blockS2] = false := by
  decide

/-! ================================================================
    NEGATIVE CONTROLS
    ================================================================ -/

/-- **Negative control (anti-strawman): every checker can return the OTHER
verdict on honest data.**  On the legitimate YES instance A = (5,5,7,5,6,6),
B = 17: `hasTP` FINDS the triple partition (so its `false` on the witness is
the output of a complete search, not a stub); the intended all-copies blocks
form a valid 12Partition; `lemmaShape` ACCEPTS both of them (so its `false`
on block1/block2 is not the output of an unsatisfiable predicate); and both
`lemma3At` and `soundnessAt` are `true` here — the refuted statements are
exactly the ones that hold on conforming data.  If any modelling definition
were degenerate (constantly false shape test, vacuous validity test, stub
search), at least one conjunct below would fail to compile. -/
theorem Algorithms_27__sol3_negative_control :
    legit Ayes 2 17 = true
  ∧ hasTP 17 2 Ayes = true
  ∧ validPartition Ayes 17 2 [confBlock1, confBlock2] = true
  ∧ lemmaShape 6 confBlock1 = true
  ∧ lemmaShape 6 confBlock2 = true
  ∧ lemma3At Ayes 17 2 [confBlock1, confBlock2] = true
  ∧ soundnessAt Ayes 17 2 [confBlock1, confBlock2] = true := by
  decide

/-- **Negative control 2: the validity checker discriminates.**  The labelled
set really has 24 distinct copies; `validPartition` rejects a duplicated
block list (copies reused), a single block (wrong count), and swapping in a
wrong-sum block; `blockSum` genuinely reads the instance values (copy (5,0)
of Aw carries value 8, and the two witness blocks sum to 68 each); `legit`
rejects an instance violating the window (4*4 = 16 < 17 and 2*9 = 18 > 17).
So `validPartition Aw 17 2 [block1, block2] = true` above is a real check of
disjointness, coverage, sizes, and sums — not a constantly-true test. -/
theorem Algorithms_27__sol3_checker_discriminates :
    (labelled 6).length = 24
  ∧ noDupP (labelled 6) = true
  ∧ validPartition Aw 17 2 [block1, block1] = false
  ∧ validPartition Aw 17 2 [block1] = false
  ∧ validPartition Aw 17 2 [block1, confBlock2] = false
  ∧ blockSum Aw [(5, 0)] = 8
  ∧ blockSum Aw block1 = 68
  ∧ blockSum Aw block2 = 68
  ∧ legit [4, 5, 5, 5, 6, 9] 2 17 = false := by
  decide

end Algorithms27sol3

\end{Vlean}
\subsection*{Proof-safety information}
The automated scan found no unfinished proofs, new axiom declarations, unsafe declarations, external code hooks, or uses of Lean's compiled decision procedure. This scan concerns proof construction only; it does not establish that the theorem matches the source claim.
\textit{Technical scan counts:} native\_decide=0, sorry=0, axiom\_decl=0, unsafe=0, extern=0. \textit{Technical tier:} KERNEL-ONLY.
\textbf{A YES means:} Lean represents the same objects and assumptions, uses the same proposed counterexample, and the final theorem statement asserts the same contradiction. Successful compilation by itself is not enough.
\clearpage
\section*{Human sign-off}
Complete this block from your own review of Questions 1--3. Do not use the automated appendix below as an answer key.\par
\textbf{Reviewer name:} \rule{9cm}{0.4pt}\par
\textbf{Date:} \rule{4cm}{0.4pt}\par
\medskip\textbf{1. Does the proof verification document accurately represent the source and challenged claim?}\par
\answerchoices\par
\decisionexplanation
\medskip\textbf{2. Does the proposed counterexample satisfy the assumptions and make that claim false?}\par
\answerchoices\par
\decisionexplanation
\medskip\textbf{3. Does the final Lean theorem statement represent and verify that same counterexample?}\par
\answerchoices\par
\decisionexplanation
\medskip\textbf{Overall decision:}\par
\checkbox\ VALIDATED (all three answers are YES)\par
\checkbox\ REJECTED (at least one answer is NO)\par
\checkbox\ NEEDS ANOTHER REVIEWER (at least one CANNOT DETERMINE and no NO)\par
\medskip\textbf{Additional comments.} Use this space for any extra source, counterexample, calculation, or Lean details that support your decisions.\par\rule{\linewidth}{0.4pt}\par\vspace{7pt}\rule{\linewidth}{0.4pt}
\clearpage
\section*{Appendix --- Prior automated review (background only)}
\textbf{These results are not an answer key.} They come from the earlier Opus/Fable verification campaign and are printed only so an auditor can see what was previously checked. The human decisions above must be made independently. Raw status codes are retained in small type for traceability.
\textbf{Overall automated disposition:} The proposed example and Lean correspondence were confirmed, but a separate search did not independently find a counterexample.
\par\textit{Audit code:} \nolinkurl{PASS_UNCORROBORATED}
\subsection*{Independent search with no access to the stored example}
\textbf{Result:} The checker could not test this claim with a finite counterexample search.
\par\textit{Audit code:} \nolinkurl{CLAIM_NOT_FINITELY_REFUTABLE}
\par\textbf{Checker's explanation:}
\begin{Vsrc}
The script did establish, exhaustively and in exact integer arithmetic, the one thing the record does license: block composition is NOT forced for a generic instance of 12Partition as the problem statement defines it (an arbitrary set of 12n positive integers). SWEEP A swept all 118755 multisets of 24 elements with values in {1..6}; (Shortened here; the complete note remains in the machine-check JSON.)
\end{Vsrc}
\par\textbf{Cases checked:}
\begin{Vsrc}
SWEEP A (exhaustive): every multiset S of 24 positive integers with values in {1,2,3,4,5,6}, i.e. all C(29,5) = 118755 multiplicity vectors (m_1..m_6) summing to 24; for each S with even total, every sub-multiset (b_1..b_6) with 0<=b_v<=m_v, sum b_v = 12 and sum v*b_v = total/2 was enumerated by exact integer recursion (n = 2, so the complement is automatically the second block), and distinct. (Shortened here; the complete note remains in the machine-check JSON.)
\end{Vsrc}
\subsection*{Independent re-check of the proposed counterexample}
\textbf{Result:} A separate checker recomputed the proposed counterexample and confirmed it.
\par\textit{Audit code:} \nolinkurl{CONFIRMED}
\subsection*{Independent review of the source-to-Lean correspondence}
\textbf{Result:} A separate reviewer judged that Lean proves the same failure using the same example.
\par\textit{Audit code:} \nolinkurl{MATCHES}
\par\textbf{Theorem summary from that reviewer:}
\begin{Vsrc}
Algorithms_27__sol3_refutation asserts, by kernel `decide`, the nine-fold conjunction on the recorded witness A = [5,5,5,5,6,8], m = 2, B = 17: (a) legit Aw 2 17 = true (|A| = 3m, sum A = mB, and the exact integer window B < 4*a and 2*a < B for every element); (b) hasTP 17 2 Aw = false, i.e. the complete anchored backtracking search finds no partition of A into triples of sum 17, so A is a NO instance of 3-Partition; (c) validPartition Aw 17 2 [block1, block2] = true, i.e. (Shortened here; the complete note remains in the machine-check JSON.)
\end{Vsrc}
\par\textbf{Difference noted by that reviewer:}
\begin{Vsrc}
None material. All three links hold and every arithmetic assertion was re-verified independently in exact integer arithmetic outside Lean (sum 34 = 2*17; window 20/20/20/20/24/32 > 17 and 10/10/10/10/12/16 < 17; exhaustive triple sums 15/16/18/19 never 17; both blocks size 12 and sum 68; multiplicities [1,0,3,4,2,2] and [3,4,1,0,2,2]; disjoint and covering all 24 labelled copies; secondary and control instances likewise). (Shortened here; the complete note remains in the machine-check JSON.)
\end{Vsrc}
\subsection*{Follow-up review}
No follow-up review was needed or recorded.
\textit{This appendix is only background. You do not need any separate machine-output files to answer the three questions.}
\end{document}